% Volume VI: Machine Learning Mathematics
% Continuity note for Chapter 35.
% Previous dependencies: Chapters 5-7 provide projections, least squares, PCA, SVD, matrix calculus, and tensor shape discipline; Chapters 18 and 31 provide information-theoretic losses, risk, empirical objectives, and inductive bias; Chapters 33-34 provide kernels, latent variables, and posterior inference.
% New durable concepts: unsupervised learning; clustering; distance and similarity; k-means objective and alternating minimization; Gaussian-mixture clustering intuition; dimensionality reduction; PCA as variance maximization and low-rank reconstruction; autoencoders; embeddings; manifold hypothesis; local neighborhoods; geodesic distance; graph Laplacian objectives; representation geometry; self-supervised targets; contrastive learning; masked prediction; InfoNCE-style losses; collapse and augmentation failures.
% Forward hooks: Chapter 37 uses representations as hidden feature maps; Chapter 38 differentiates autoencoder and contrastive objectives; Chapter 42 studies representation geometry and robustness; Chapter 47-48 connect neural coding and latent neural dynamics to manifolds; Chapter 53 returns to PCA-to-neural-manifold case studies.
% Notation commitments: Unlabeled data are x_i in R^d; data matrices are centered when stated; cluster assignments are c_i in {1,...,K}; centroids are mu_k; low-dimensional coordinates are z_i in R^r; encoders are e_phi, decoders are d_theta; representation maps are h or h_theta; neighborhood graphs have weights W_{ij} and Laplacian L=D-W.
\section{Chapter 35: Unsupervised Learning, Manifolds, and Representation}
\label{ch:unsupervised-learning-manifolds-representation}

Supervised learning starts with input-output pairs. Much of neuroscience and artificial intelligence starts with something harder: a collection of observations without labels that say what each observation means. A population recording may contain thousands of neural activity vectors during behavior. A video stream contains images before any class label is attached. A language model sees sequences of tokens and must create internal features before a downstream task is specified.

The central problem is to discover useful structure in data. Useful does not mean uniquely correct. Clusters depend on a distance, dimensionality reduction depends on what variation is preserved, manifold learning depends on local neighborhoods, and self-supervised learning depends on the prediction task used to create training signals from the data itself. This chapter teaches the mathematical objects behind those choices.

\paragraph{Chapter problem.}
How can unlabeled data be organized into clusters, low-dimensional coordinates, manifolds, and learned representations that preserve structure useful for neuroscience and AI?

\paragraph{Section map.}
Section 35.1 introduces clustering through distances, assignments, centroids, and mixture-model intuition. Section 35.2 develops dimensionality reduction through PCA, low-rank reconstruction, embeddings, and autoencoders. Section 35.3 explains the manifold hypothesis, neighborhood graphs, geodesic distance, and Laplacian objectives. Section 35.4 previews self-supervised learning through pretext tasks, contrastive objectives, masked prediction, and representation collapse.

\subsection{35.1 Clustering}

\paragraph{Guiding question.}
Given unlabeled points, when is it meaningful to divide them into groups, and what objective is being optimized?

\paragraph{Beginner-facing intuition.}
Clustering tries to replace a cloud of points by a small number of groups. A cluster should contain points that are similar to each other and different from points in other clusters. The word ``similar'' is not automatic. It may mean small Euclidean distance between feature vectors, high correlation between neural response profiles, shared graph neighborhoods, or high probability under the same latent component.

The simplest model is centroid clustering. Each cluster has a representative point, its centroid. Each data point is assigned to the closest centroid. Then each centroid is updated to the average of the points assigned to it. Repeating these two steps gives the \(k\)-means algorithm.

\paragraph{Formal definitions.}
Let
\[
x_1,\ldots,x_n\in\mathbb R^d
\]
be unlabeled data. A hard \(K\)-clustering assigns each point a label
\[
c_i\in\{1,\ldots,K\}.
\]
For \(k\)-means, each cluster has a centroid
\[
\mu_k\in\mathbb R^d.
\]
The \(k\)-means objective is
\[
\boxed{
J(c,\mu)=\sum_{i=1}^n \|x_i-\mu_{c_i}\|_2^2.
}
\]
The goal is to choose assignments \(c=(c_1,\ldots,c_n)\) and centroids \(\mu=(\mu_1,\ldots,\mu_K)\) that make within-cluster squared distances small.

A probabilistic alternative is a Gaussian mixture model. It introduces a latent cluster variable \(Z_i\in\{1,\ldots,K\}\) and parameters
\[
P(Z_i=k)=\pi_k,\qquad X_i\mid Z_i=k\sim\mathcal N(m_k,\Sigma_k).
\]
Then cluster membership is inferred from posterior responsibilities \(P(Z_i=k\mid x_i)\), connecting this section to Chapter 34.

\paragraph{Core mechanism: the two \(k\)-means updates decrease the objective.}
Fix the centroids \(\mu_1,\ldots,\mu_K\). For each point \(x_i\), the objective term involving its assignment is
\[
\|x_i-\mu_{c_i}\|^2.
\]
Choosing
\[
c_i\in\operatorname*{arg\,min}_{1\leq k\leq K}\|x_i-\mu_k\|^2
\]
minimizes that term independently of the other assignments, so the assignment step cannot increase \(J\).

Now fix the assignments. For one cluster \(C_k=\{i:c_i=k\}\), its contribution is
\[
J_k(\mu_k)=\sum_{i\in C_k}\|x_i-\mu_k\|^2.
\]
Differentiate with respect to \(\mu_k\):
\[
\nabla_{\mu_k}J_k=2\sum_{i\in C_k}(\mu_k-x_i).
\]
Setting the gradient to zero gives
\[
\boxed{
\mu_k=\frac{1}{|C_k|}\sum_{i\in C_k}x_i.
}
\]
Thus the centroid update also minimizes \(J\) with the other variables fixed. Alternating these steps monotonically decreases the objective until a local optimum or fixed point is reached. The problem is not convex jointly in assignments and centroids, so different initializations can give different clusterings.

\paragraph{Concrete example.}
Take four points on the line:
\[
x_1=0,\quad x_2=1,\quad x_3=5,\quad x_4=6,
\]
and \(K=2\). Initialize \(\mu_1=0\), \(\mu_2=6\). The assignment step gives \(C_1=\{0,1\}\) and \(C_2=\{5,6\}\). The update step gives
\[
\mu_1=(0+1)/2=0.5,\qquad \mu_2=(5+6)/2=5.5.
\]
The within-cluster objective is
\[
(0-0.5)^2+(1-0.5)^2+(5-5.5)^2+(6-5.5)^2=1.
\]
This clustering matches the visual gap between \(1\) and \(5\). If the data were stretched or rescaled feature by feature, the answer could change.

\paragraph{Computational neuroscience example.}
Suppose \(x_i\in\mathbb R^N\) is the population spike-count vector from trial \(i\) across \(N\) neurons. Clustering might reveal trials with different behavioral states, attention levels, or stimulus conditions. The mathematical output is a partition of response vectors. It is not by itself proof that the brain has exactly \(K\) discrete internal states. Validation requires held-out behavior, stimulus labels, perturbations, or a generative model.

\paragraph{AI and machine-learning example.}
In image or embedding databases, clustering groups examples whose learned feature vectors are close. Product quantization and vector search use clustering to compress large embedding collections. Gaussian mixtures provide soft cluster probabilities, which can be useful when an example lies between groups.

\paragraph{Common misconceptions and failure modes.}
Clusters are not intrinsic unless the metric, representation, and number of clusters are justified. \(k\)-means prefers roughly spherical clusters of comparable scale because it uses squared Euclidean distance to centroids. Empty clusters can occur during updates. Local minima make initialization important. A low within-cluster objective does not mean the clusters correspond to meaningful biological or semantic categories.

\paragraph{Exercises.}
\begin{enumerate}[label=\textbf{35.1.\arabic*.}, leftmargin=*]
\item \textbf{Mechanical.} For \(x_1=(0,0)\), \(x_2=(1,0)\), \(x_3=(5,0)\), centroids \(\mu_1=(0,0)\), \(\mu_2=(4,0)\), assign each point to its nearest centroid.
\item \textbf{Proof or derivation.} Derive the centroid update \(\mu_k=|C_k|^{-1}\sum_{i\in C_k}x_i\) by differentiating the within-cluster squared-error objective.
\item \textbf{Numerical/computational.} Run one full \(k\)-means assignment/update step for the four one-dimensional points \(0,2,8,10\) initialized at centroids \(0\) and \(10\).
\item \textbf{Interpretation.} In a neural trial-clustering analysis, what extra evidence would support interpreting clusters as behavioral states rather than recording noise?
\item \textbf{Misconception/debugging.} A student says, ``The clustering algorithm found three clusters, so the data truly have three natural categories.'' Explain why this conclusion depends on \(K\), distance, representation, and validation.
\end{enumerate}

\subsection{35.2 Dimensionality reduction}

\paragraph{Guiding question.}
How can high-dimensional data be represented by fewer coordinates while preserving the structure we care about?

\paragraph{Beginner-facing intuition.}
High-dimensional observations often vary along fewer effective directions than their ambient dimension suggests. A neural population may contain hundreds of recorded neurons, but task-related activity may lie near a low-dimensional subspace. An image has thousands of pixels, but much variation may be described by lighting, pose, or object identity. Dimensionality reduction replaces each high-dimensional point \(x_i\in\mathbb R^d\) by a lower-dimensional coordinate \(z_i\in\mathbb R^r\), with \(r<d\).

The simplest method is principal component analysis, or PCA. PCA finds orthogonal directions of largest variance in centered data. Keeping the top \(r\) directions gives a linear low-dimensional representation and a low-rank reconstruction.

\paragraph{Formal definitions.}
Let \(X\in\mathbb R^{n\times d}\) be a centered data matrix whose \(i\)-th row is \(x_i^\top\), so \(\sum_i x_i=0\). A rank-\(r\) linear representation chooses a matrix
\[
U_r\in\mathbb R^{d\times r},\qquad U_r^\top U_r=I_r,
\]
and represents each point by
\[
z_i=U_r^\top x_i\in\mathbb R^r.
\]
The reconstruction is
\[
\widehat x_i=U_r z_i=U_rU_r^\top x_i.
\]
PCA chooses \(U_r\) to maximize projected variance
\[
\operatorname{tr}(U_r^\top S U_r),
\qquad
S=\frac1n X^\top X,
\]
or equivalently to minimize total squared reconstruction error
\[
\sum_{i=1}^n \|x_i-U_rU_r^\top x_i\|^2.
\]

An autoencoder generalizes this idea by learning an encoder and decoder
\[
e_\phi:\mathbb R^d\to\mathbb R^r,\qquad d_\theta:\mathbb R^r\to\mathbb R^d
\]
and minimizing a reconstruction loss such as
\[
\sum_{i=1}^n \|x_i-d_\theta(e_\phi(x_i))\|^2.
\]

\paragraph{Core theorem: PCA solves the best linear reconstruction problem.}
Let \(S=(1/n)X^\top X\) have eigenvalues
\[
\lambda_1\geq\lambda_2\geq\cdots\geq\lambda_d\geq0
\]
with orthonormal eigenvectors \(u_1,\ldots,u_d\). The rank-\(r\) PCA subspace is
\[
U_r=[u_1\ \cdots\ u_r].
\]
It maximizes \(\operatorname{tr}(U^\top S U)\) over all \(U^\top U=I_r\) and minimizes squared reconstruction error among all \(r\)-dimensional orthogonal projections.

\emph{Derivation.} For any orthonormal \(U\),
\[
\sum_i \|x_i-UU^\top x_i\|^2
=\sum_i \|x_i\|^2-\sum_i \|U^\top x_i\|^2
=\operatorname{tr}(X^\top X)-n\,\operatorname{tr}(U^\top S U).
\]
The first term is fixed, so minimizing reconstruction error is the same as maximizing projected variance. Expanding \(U\) in the eigenbasis of \(S\) shows that the maximum is obtained by placing the \(r\) orthonormal columns along the eigenvectors with largest eigenvalues. The maximum captured variance is \(\sum_{j=1}^r\lambda_j\).

\paragraph{Concrete example.}
Consider centered two-dimensional data with covariance
\[
S=\begin{pmatrix}4&0\\0&1\end{pmatrix}.
\]
The eigenvectors are \(e_1=(1,0)^\top\) and \(e_2=(0,1)^\top\), with eigenvalues \(4\) and \(1\). One-dimensional PCA chooses \(U_1=e_1\). The coordinate of \(x=(3,2)\) is
\[
z=e_1^\top x=3,
\]
and the reconstruction is
\[
\widehat x=e_1z=(3,0).
\]
The second coordinate is discarded because it lies along the lower-variance direction.

\paragraph{Computational neuroscience example.}
Let \(X\in\mathbb R^{T\times N}\) contain the trial-averaged activity of \(N\) neurons over \(T\) time bins during a reaching task. PCA finds population activity directions that explain large variance over time. Plotting the first two or three scores can reveal a low-dimensional neural trajectory. This is a descriptive geometry of activity, not proof that the nervous system explicitly computes PCA.

\paragraph{AI and machine-learning example.}
Autoencoders learn nonlinear dimensionality reduction by training an encoder-decoder pair with backpropagation. A bottleneck code \(z=e_\phi(x)\) can compress images, denoise inputs, or initialize representations for later tasks. Unlike PCA, a nonlinear autoencoder can represent curved structure, but it can also memorize or learn a code that is hard to interpret.

\paragraph{Common misconceptions and failure modes.}
Low-dimensional does not always mean interpretable. PCA directions maximize variance, not task relevance. Features with larger units can dominate PCA unless data are scaled appropriately. A two-dimensional visualization may hide important structure in later components. Autoencoder reconstruction quality does not guarantee that the latent coordinates are useful for classification, control, or neuroscience interpretation.

\paragraph{Exercises.}
\begin{enumerate}[label=\textbf{35.2.\arabic*.}, leftmargin=*]
\item \textbf{Mechanical.} If \(X\in\mathbb R^{100\times 50}\) is centered and \(U\in\mathbb R^{50\times 3}\), state the shapes of \(Z=XU\) and \(\widehat X=Z U^\top\).
\item \textbf{Proof or derivation.} Starting from \(\|x-UU^\top x\|^2\), show that reconstruction error equals \(\|x\|^2-\|U^\top x\|^2\) when \(U^\top U=I\).
\item \textbf{Numerical/computational.} For covariance \(\operatorname{diag}(9,4,1)\), compute the fraction of variance explained by the first two principal components.
\item \textbf{Interpretation.} Why might the first principal component of neural data reflect movement speed, arousal, or recording drift rather than the task variable of interest?
\item \textbf{Misconception/debugging.} A student says, ``PCA found the true latent variables.'' Explain why PCA finds variance-maximizing linear directions, not necessarily causal or mechanistic latent variables.
\end{enumerate}

\subsection{35.3 Manifold learning}

\paragraph{Guiding question.}
What changes when data lie near a curved low-dimensional set rather than a flat linear subspace?

\paragraph{Beginner-facing intuition.}
A line is a one-dimensional subspace only if it is straight and passes through the origin. A circle is also one-dimensional locally, but it is curved and cannot be represented well by one global linear coordinate without cutting it. Manifold learning starts from the idea that high-dimensional data may lie near a low-dimensional curved surface. Locally the data look simple; globally the geometry may bend.

The practical problem is that only samples are observed. Manifold methods approximate local geometry using neighborhoods, then build global coordinates that preserve neighborhood, distance, or smoothness structure.

\paragraph{Formal definitions.}
A \(r\)-dimensional manifold \(\mathcal M\subseteq\mathbb R^d\) is a set that locally resembles \(\mathbb R^r\). Informally, near each point \(x\in\mathcal M\), there are coordinates \(z\in\mathbb R^r\) that describe nearby points. The \emph{manifold hypothesis} says that observed high-dimensional data concentrate near such a low-dimensional set.

Given sample points \(x_1,\ldots,x_n\), a neighborhood graph has weights
\[
W_{ij}\geq0
\]
that are large when points \(i\) and \(j\) are neighbors. The degree matrix is
\[
D_{ii}=\sum_j W_{ij},
\]
and the graph Laplacian is
\[
L=D-W.
\]
An embedding assigns coordinates \(y_i\in\mathbb R^r\) to each data point. Laplacian eigenmaps choose coordinates that keep neighbors close by minimizing
\[
\boxed{
\frac12\sum_{i,j}W_{ij}\|y_i-y_j\|^2
=\operatorname{tr}(Y^\top L Y),
}
\]
under normalization constraints that prevent the trivial solution \(Y=0\).

\paragraph{Core mechanism: the graph Laplacian penalizes nonsmooth embeddings.}
Let \(y\in\mathbb R^n\) assign one scalar coordinate \(y_i\) to each data point. Then
\[
y^\top L y
=y^\top(D-W)y
=\sum_i D_{ii}y_i^2-\sum_{i,j}W_{ij}y_i y_j.
\]
Because \(D_{ii}=\sum_jW_{ij}\) and \(W\) is symmetric,
\[
\boxed{
y^\top L y
=\frac12\sum_{i,j}W_{ij}(y_i-y_j)^2.
}
\]
Thus the objective is small when neighboring points in the data graph receive similar embedding coordinates. The nontrivial eigenvectors of \(L\), or of a normalized Laplacian, provide smooth coordinates over the graph.

\paragraph{Concrete example.}
For three points in a chain \(1-2-3\) with weights \(W_{12}=W_{21}=1\), \(W_{23}=W_{32}=1\), and no edge \(1-3\),
\[
L=\begin{pmatrix}
1&-1&0\\
-1&2&-1\\
0&-1&1
\end{pmatrix}.
\]
For coordinate values \(y=(0,1,2)^\top\),
\[
y^\top L y=(0-1)^2+(1-2)^2=2.
\]
For \(y=(0,2,0)^\top\),
\[
y^\top L y=(0-2)^2+(2-0)^2=8.
\]
The second assignment changes sharply across neighboring edges, so it is less smooth on the graph.

\paragraph{Computational neuroscience example.}
During a navigation task, a neural population vector may trace a ring-like or torus-like manifold as head direction, position, or phase changes. Euclidean distance between two high-dimensional activity vectors may be misleading if the true latent variable is circular. Neighborhood-based manifold learning can reveal the intrinsic geometry, but the result depends on sampling density, noise, and the chosen neighborhood scale.

\paragraph{AI and machine-learning example.}
UMAP, diffusion maps, Isomap, and related methods build graphs from nearest neighbors and compute embeddings that preserve local or geodesic structure. Deep representation learning can be seen as learning a map \(h_\theta(x)\) in which the data manifold is easier to separate, interpolate, or model.

\paragraph{Common misconceptions and failure modes.}
The manifold hypothesis is an assumption, not a theorem about every dataset. Neighborhood graphs can connect points that are close in ambient space but far along the manifold, especially in noisy or sparsely sampled regions. Two-dimensional embeddings are visual summaries, not complete proofs of cluster or manifold structure. Geodesic distance along a manifold differs from Euclidean chord distance, but estimating it from data is fragile.

\paragraph{Exercises.}
\begin{enumerate}[label=\textbf{35.3.\arabic*.}, leftmargin=*]
\item \textbf{Mechanical.} For a graph with weights \(W_{12}=W_{21}=2\), \(W_{23}=W_{32}=1\), write the degree matrix \(D\) and Laplacian \(L\).
\item \textbf{Proof or derivation.} Derive \(y^\top L y=\frac12\sum_{i,j}W_{ij}(y_i-y_j)^2\) for a symmetric weight matrix.
\item \textbf{Numerical/computational.} For the three-node chain in the section, compute \(y^\top L y\) for \(y=(1,1,1)^\top\) and interpret the result.
\item \textbf{Interpretation.} Why can a Swiss-roll dataset require a nonlinear embedding even though it is locally two-dimensional?
\item \textbf{Misconception/debugging.} A student says, ``A UMAP plot with separated islands proves there are discrete neural states.'' Explain why neighborhood choice, sampling, and visualization distortion must be checked.
\end{enumerate}

\subsection{35.4 Self-supervised learning preview}

\paragraph{Guiding question.}
How can a system learn useful representations from data without externally supplied labels?

\paragraph{Beginner-facing intuition.}
Self-supervised learning creates training signals from the data itself. Instead of asking for a human label, the model predicts a hidden part of the input, a future observation, another view of the same example, or whether two views belong together. The goal is not merely to solve the pretext task; it is to learn a representation \(h_\theta(x)\) that transfers to later tasks.

This idea is important for neuroscience because sensory systems receive enormous streams of structured input without explicit labels for every moment. It is also central in modern AI, where language models, vision models, and audio models learn from prediction problems constructed from raw data.

\paragraph{Formal definitions.}
Let \(X\sim P_X\) be unlabeled data. A representation map is
\[
h_\theta:\mathcal X\to\mathbb R^m.
\]
A self-supervised task defines a loss using transformations, masks, or future observations derived from \(X\). In masked prediction, an input is split into observed and hidden parts,
\[
X=(X_{\mathrm{obs}},X_{\mathrm{mask}}),
\]
and the model predicts \(X_{\mathrm{mask}}\) from \(X_{\mathrm{obs}}\). In contrastive learning, two augmented views \(X^+\) and \(\widetilde X^+\) of the same example form a positive pair, while views from other examples are negatives.

One common contrastive loss for an anchor representation \(q\), a positive key \(k^+\), and negative keys \(k_1^-,\ldots,k_M^-\) is
\[
\boxed{
\ell_{\mathrm{NCE}}
=-\log
\frac{\exp(\operatorname{sim}(q,k^+)/\tau)}
{\exp(\operatorname{sim}(q,k^+)/\tau)+\sum_{m=1}^M\exp(\operatorname{sim}(q,k_m^-)/\tau)}.
}
\]
Here \(\operatorname{sim}\) is a similarity score such as a dot product or cosine similarity, and \(\tau>0\) is a temperature.

\paragraph{Core mechanism: contrastive learning is classification among candidates.}
Define scores
\[
s_0=\operatorname{sim}(q,k^+)/\tau,\qquad
s_m=\operatorname{sim}(q,k_m^-)/\tau.
\]
The softmax probability assigned to the positive item is
\[
p_0=\frac{e^{s_0}}{\sum_{j=0}^M e^{s_j}}.
\]
The contrastive loss is
\[
\ell_{\mathrm{NCE}}=-\log p_0.
\]
Thus the model is trained to classify which candidate key is the matching view. If \(s_0\) increases while other scores stay fixed, \(p_0\) increases and the loss decreases. If a negative is too similar to the anchor, it competes in the denominator and increases the loss. This is the same cross-entropy mechanism introduced in Chapter 18, now applied to representation pairs.

\paragraph{Concrete example.}
Suppose one positive and two negatives have scores
\[
s_0=2,\qquad s_1=1,\qquad s_2=0.
\]
Then
\[
p_0=\frac{e^2}{e^2+e^1+e^0}\approx \frac{7.39}{11.11}\approx0.665,
\]
and
\[
\ell=-\log(0.665)\approx0.408.
\]
If the positive score falls to \(0.5\) while negatives stay at \(1\) and \(0\), the probability assigned to the positive decreases, and the loss increases.

\paragraph{Computational neuroscience example.}
A predictive-coding-inspired self-supervised task might predict a future sensory frame, a masked segment of a sound, or the next neural population state from recent history. The mathematical claim is that prediction can shape representations. It does not imply that every cortical circuit literally optimizes a particular contrastive or masked-token objective.

\paragraph{AI and machine-learning example.}
Language models predict masked or next tokens from context. Vision systems learn by making two augmentations of the same image have nearby representations while pushing apart unrelated images, or by predicting masked image patches. Speech models may predict masked acoustic features. These methods can produce features useful for classification, retrieval, transfer learning, and control.

\paragraph{Common misconceptions and failure modes.}
Self-supervised does not mean objective-free; the constructed task is a strong inductive bias. Bad augmentations can remove information needed later or make shortcuts possible. Contrastive learning can treat semantically similar examples as negatives. Noncontrastive methods must avoid representation collapse, where \(h_\theta(x)\) becomes nearly constant. High pretext-task accuracy does not guarantee useful downstream representations.

\paragraph{Exercises.}
\begin{enumerate}[label=\textbf{35.4.\arabic*.}, leftmargin=*]
\item \textbf{Mechanical.} In the contrastive loss, identify the anchor, positive key, negative keys, similarity function, and temperature.
\item \textbf{Proof or derivation.} Show that \(\ell_{\mathrm{NCE}}=-\log p_0\) is the ordinary cross-entropy loss for classifying the positive key among \(M+1\) candidates.
\item \textbf{Numerical/computational.} Compute the contrastive loss when scores are \(s_0=1\), \(s_1=0\), \(s_2=0\). Use natural logarithms.
\item \textbf{Interpretation.} Give one example of a self-supervised task for neural recordings and state what representation you would hope it learns.
\item \textbf{Misconception/debugging.} A student says, ``No labels means no assumptions.'' Explain how masking rules, augmentations, prediction horizons, and negative sampling encode assumptions.
\end{enumerate}

Unsupervised learning is not one method but a family of ways to expose structure without external labels. Clustering gives groups, PCA and autoencoders give coordinates, manifold learning gives local-to-global geometry, and self-supervised learning gives trainable objectives for representation maps. The next chapter adds action and reward, where representations and states must support decisions over time.

\clearpage
